\title{Porting a continuous-energy Monte Carlo to CUDA in pure Rust:\\[0.3em]
\normalsize What we built, what it weighs, and what we
believed that wasn't true}
\author{F\'abio Rodrigues\\\small \texttt{sorcerer86pt@gmail.com}}
\date{2026}

\maketitle

\begin{abstract}
We report the engineering, observability, and honest performance
findings of porting \texttt{open\_rust\_mc} --- a pure-Rust
continuous-energy neutron / photon Monte Carlo transport code with
recursive-CSG geometry, SVD-compressed cross sections, and ICSBEP
regression substrate --- to NVIDIA CUDA via the \texttt{cudarc}
crate and NVRTC. The CPU and GPU backends share a single
\texttt{EigenvalueRunner} trait so any benchmark runs through either
path with the same scene JSON, the same nuclear data, and the same
acceptance envelope.

The paper documents three engineering decisions that turned out to
matter more than the kernel-level optimisation we initially
prioritised: \textbf{(i)} a process-wide per-nuclide device-buffer
cache keyed on \texttt{Arc::as\_ptr} that amortises 50\,$\times$
+ 50~MB-class HtoD copies across a multi-case sweep; \textbf{(ii)}
a parallel cache for S($\alpha$,$\beta$) thermal-scattering payloads
whose absence quietly wasted tens of seconds per case on benchmarks
that share moderators (every PWR pin cell using
\texttt{c\_H\_in\_H2O}, every Be-reflected fast-metal case using
\texttt{c\_Be}); \textbf{(iii)} adaptive NVRTC arch detection so a
single binary targets RTX A1000 / RTX 3080 / A100 / H100 / RTX 5090
without rebuilds.

Three findings we did not expect: GPU SVD reconstruction, despite a
$2.6$--$2.8\,\times$ \emph{per-XS-lookup} micro-speedup over a
single-threaded CPU loop, is $1.3\,\times$ \emph{slower} than the
20-core rayon-parallel CPU SVD at full transport integration on
Godiva (and slower still vs.\ GPU pointwise table on PWR until a
synthesised MT$=4$ + per-level CDF closed the gap). The ``GPU
recursive transport 6.74$\,\times$ CPU'' headline that propagated
through earlier sessions is a microbenchmark of constant-XS
geometry traversal, not a measurement of integrated transport.
Refill-pool occupancy gains (PHYSOR 2022 Optimization~F) only
appear in the mid-curve regime where the device is under-occupied
but not yet saturated; both extremes give no measurable benefit.

We outline a structured observability path
(\texttt{-{}-debug-metrics}, exposed via a new
\texttt{gpu\_debug\_metrics()} PyO3 entry point that snapshots cache
entries, hit counts, byte usage, and the detected
compute-capability) that lets a sweep harness verify the cache
machinery actually fires at runtime rather than relying on
engineering belief.

\textbf{Code:}
\href{https://github.com/sorcerer86pt/open_rust_mc}%
{github.com/sorcerer86pt/open\_rust\_mc}
\end{abstract}
